\clearpage\thispagestyle{empty}
\pdfbookmark[0]{Quellen zur Besprechung von Hardings Himmelsatlas}{harding1812-sources}
\begin{center}\textbf{Quellen zur Besprechung von Hardings Himmelsatlas}\end{center}
\noindent Diese KI-Übersetzung ins heutige Deutsch folgt Carl Friedrich Gauß, \textit{Werke}, Bd.~VI (Göttingen, 1874), S.~541--542. Die Besprechung erschien zuerst unter Göttingen in den \textit{Göttingischen gelehrten Anzeigen}, Stück~158, 3. Oktober 1812, S.~1569--1571.
\par\smallskip\noindent Im Text der Werkausgabe fehlt das Komma zwischen Skorpion und Schütze, das im Erstdruck steht. Die Schreibweisen \textit{nordlich} und \textit{Verschliess} sind beibehalten. \textit{Verschliess} bezeichnet hier den Verkauf der Karten. Hardings späterer \href{https://www.davidrumsey.com/luna/servlet/detail/RUMSEY~8~1~372238~90138977:Title--Atlas-novus-coelestis-XXVII-}{Atlas von 1822} ist in der David Rumsey Historical Map Collection digital zugänglich.
\par\smallskip\noindent \href{https://gdz.sub.uni-goettingen.de/id/PPN236007467}{Text der Werkausgabe}. Erstdruck, \href{https://images.p.sub.uni-goettingen.de/iiif/image/gdz:PPN319721507_1812:00001613/full/1650,/0/default.jpg}{S.~1569}, \href{https://images.p.sub.uni-goettingen.de/iiif/image/gdz:PPN319721507_1812:00001614/full/1650,/0/default.jpg}{S.~1570} und \href{https://images.p.sub.uni-goettingen.de/iiif/image/gdz:PPN319721507_1812:00001615/full/1650,/0/default.jpg}{S.~1571}.
\par\smallskip\noindent Übersetzung und Satz wurden von KI erstellt.
\clearpage
